\section{Methods}
\label{sec:methods}

Three choices in this section are conventions rather than results, and each has cost this
project a wrong number at least once. They are the coordinate in which a spectrum is taken,
the power of the amplitude on which a sparsity statistic is computed, and what a random
control is allowed to randomise. We state each with its formula and its parameter, in the manner of
\citet{doshi2024grokking}'s footnote, because the literature does not agree on them.

\subsection{Model and training}
\label{sec:setup-training}

% STE descriptive
All runs use one transformer architecture from \citet{nanda2023progress}. The model has a
single layer, $d_{\text{model}} = 128$, four attention heads of width $32$, and a ReLU MLP of
width $512$. It has no layer normalisation or bias terms. The logits come from the final
position as $(z + \mathrm{ReLU}(z W_{\text{in}}) W_{\text{out}}) W_{U}$, where $z$ is the
post-attention residual stream. A learned positional embedding $W_{\text{pos}}$ adds to the
token embedding $W_{E}$, and attention uses a causal mask. Outside the sweep of
Section~\ref{sec:causal-init}, each weight matrix starts from a zero-mean Gaussian with
standard deviation $1/\sqrt{512}$ for $W_{\text{out}}$ and $1/\sqrt{128}$ for all other
matrices.

Training is full-batch AdamW with learning rate $10^{-3}$, weight decay $1.0$,
$\beta = (0.9, 0.98)$ and $\epsilon = 10^{-8}$. The loss is cross-entropy on the $30\%$
training split, for $40{,}000$ steps. The run seed sets both the initial weights and the
random split. PyTorch logs test accuracy and both losses every $200$ steps, and the engine
logs them every $100$ steps. Test accuracy is measured on the pairs held out from
training. The grokking step of a run is the first logged step at which test accuracy is above
$0.99$.

The weight decay is the setting under which grokking is reliably observed on this task
family \citep{power2022grokking}. It is held fixed across every modulus so that the modulus
is the only variable.

Two independent implementations produce the results in this paper. The breadth, $23$ moduli
with up to five seeds, comes from a PyTorch implementation. The same architecture is also
implemented on an autograd engine written from scratch for this project. The engine
reproduces the PyTorch forward pass and gradients to machine precision and is used as an
independent check at four moduli. Section~\ref{sec:impl} treats the two arms and the one
systematic difference between them.
% STE end

\subsection{Coordinates}
\label{sec:methods-coords}
\label{sec:setup-coords}

Two coordinate systems appear throughout, and every character index in this paper names which
one it is in. Where $(\mathbb{Z}/q\mathbb{Z})^{\times}$ is cyclic we fix a primitive root $g$
and use \emph{discrete-log coordinates}, $x = g^{\,\ell}$, in which multiplication is addition
of $\ell$. Where it is not cyclic, as at $n = 119$, $n = 120$, $n = 128$ and in every stratum
whose local group is a product, there is no discrete log. There we use \emph{exponent-tuple
coordinates}, given by a fixed decomposition into cyclic factors. The choice of $g$ is a choice
of labelling and not of structure; Proposition~\ref{prop:generator} records that
the readout is exactly invariant to it.

All spectral statistics on a multiplication run are taken in a coordinate in which
multiplication is addition. Where $(\mathbb{Z}/q\mathbb{Z})^{\times}$ is cyclic of order
$\varphi(q)$, a primitive root $g$ gives the discrete-log isomorphism
\begin{equation}
  \label{eq:dlog}
  \mathrm{dlog}_{g} : (\mathbb{Z}/q\mathbb{Z})^{\times} \;\xrightarrow{\ \sim\ }\;
  \mathbb{Z}/\varphi(q)\mathbb{Z},
  \qquad g^{\ell} \longmapsto \ell,
  \qquad \mathrm{dlog}_{g}(xy) = \mathrm{dlog}_{g}(x) + \mathrm{dlog}_{g}(y) ,
\end{equation}
and the characters of the unit group are
\begin{equation}
  \label{eq:character}
  \chi_{k}(x) \;=\; \exp\!\Bigl( \tfrac{2\pi i\, k\, \mathrm{dlog}_{g}(x)}{\varphi(q)} \Bigr),
  \qquad k \in \mathbb{Z}/\varphi(q)\mathbb{Z} .
\end{equation}
The frequency index $k$ lives in $\mathbb{Z}/\varphi(q)\mathbb{Z}$, not in
$\mathbb{Z}/n\mathbb{Z}$, and the two are routinely confused. Where the unit group is not
cyclic, the structure theorem gives $G \cong \mathbb{Z}_{o_{1}} \times \cdots \times
\mathbb{Z}_{o_{r}}$ with generators $g_{1}, \dots, g_{r}$; every unit is uniquely
$x = \prod_{j} g_{j}^{\,e_{j}}$, and the \emph{product characters} are
\begin{equation}
  \label{eq:prodchar}
  \chi_{\mathbf{k}}(x) \;=\;
  \exp\!\Bigl( 2\pi i \sum_{j=1}^{r} \tfrac{k_{j} e_{j}}{o_{j}} \Bigr),
  \qquad \mathbf{k} \in \widehat{G} \cong \prod_{j} \mathbb{Z}_{o_{j}},
  \qquad |\widehat{G}| = |G| .
\end{equation}
The logit grid is then a $2r$-dimensional array. In both cases $\chi_{\mathbf{k}}(xy) =
\chi_{\mathbf{k}}(x)\chi_{\mathbf{k}}(y)$, so multiplication becomes addition of exponent
tuples, which is what makes a clock expressible at all. For a matrix $M \in
\mathbb{R}^{G \times D}$ indexed by units, the character transform is
\begin{equation}
  \label{eq:transform}
  \widehat{M}_{\mathbf{k}} \;=\; \sum_{x \in G} \overline{\chi_{\mathbf{k}}(x)}\, M_{x}
  \;\in\; \mathbb{C}^{D} .
\end{equation}

A statistic reported in raw residue coordinates on a multiplication run is not a weaker
measurement, it is a different one: the same model's top eight logit components read $17.0\%$ ``$a+b$'' in raw coordinates
and $98.4\%$ in discrete-log coordinates.

\paragraph{Invariance to the generator.} The primitive root $g$ in \eqref{eq:dlog} is
arbitrary (our implementation returns the smallest), so nothing claimed about a circuit
may depend on it.

\begin{proposition}[Generator equivariance]
\label{prop:generator}
Let $\varphi = \varphi(q)$ and let $t$ satisfy $\gcd(t, \varphi) = 1$. Replacing $g$ by $g^{t}$
sends $\mathrm{dlog}_{g} \mapsto t^{-1} \mathrm{dlog}_{g} \bmod \varphi$ and therefore permutes
the character index by $k \mapsto t k \bmod \varphi$. Consequently
\begin{equation}
  \label{eq:equivariance}
  \mathrm{Gini},\;\; \mathrm{PR},\;\;
  \mathcal{L}(L),\;\; \mathcal{L}(\Pi_{K} L),\;\; \mathcal{L}(\Pi_{K^{c}} L),\;\; p
  \quad\text{are invariant,}
  \qquad
  K \;\longmapsto\; tK \bmod \varphi .
\end{equation}
\end{proposition}

\begin{proof}
$k \mapsto tk$ is a bijection of $\mathbb{Z}/\varphi\mathbb{Z}$, so it permutes the amplitude
multiset $\{A_{k}\}$; \eqref{eq:gini} and \eqref{eq:pr} are symmetric functions of that
multiset. It also maps the diagonal to itself, $(k,k) \mapsto (tk, tk)$, so $\Pi_{K}$ and
$\Pi_{K^{c}}$ in \eqref{eq:ablation} are conjugated by a permutation and their losses are
unchanged. The $p$ of \eqref{eq:equivariance} is the exact permutation $p$, a function of those
losses and of the uniform distribution over sets of fixed cardinality, which the permutation
preserves. Its Monte-Carlo estimate $\hat{p}$ of \eqref{eq:permp} is therefore invariant in
distribution, not draw by draw, because a fixed random seed selects different characters
under a different generator.
\end{proof}

Gini, PR and the three losses are verified numerically to $10^{-10}$ at $113$, $121$ and $125$,
and the losses per cyclic component at the non-cyclic moduli. For $\hat{p}$ the check is that
the $p < 0.01$ verdict does not change. A key-frequency value is therefore not a property of a model. Only
the set up to multiplication by a unit modulo $\varphi$ is, and this paper states key sets in
that sense. Writing ``key frequencies $\{11, 29, 33, 49\}$ at $n = 121$'' names a representative of an orbit, not
an invariant.

\subsection{Sparsity statistics, and the three conventions in the literature}
\label{sec:methods-sparsity}

Let $M$ be an embedding or activation matrix indexed by residues. For frequency $k$ we form
the sine and cosine projections $s_{k}, c_{k}$ and combine them into a single
\emph{amplitude}
\begin{equation}
  \label{eq:amplitude}
  A_{k} \;=\; \sqrt{\lVert s_{k} \rVert^{2} + \lVert c_{k} \rVert^{2}},
  \qquad k = 1, \dots, \lfloor \varphi/2 \rfloor,
\end{equation}
dropping the DC component $k = 0$. Two statistics are computed from $\{A_{k}\}$, and they are
computed on different powers of it.

Gini is computed on the amplitude. With $A_{(1)} \le \cdots \le A_{(M)}$ the sorted
amplitudes,
\begin{equation}
  \label{eq:gini}
  \mathrm{Gini}(A) \;=\;
  \frac{\sum_{i=1}^{M} (2i - M - 1)\, A_{(i)}}{M \sum_{i=1}^{M} A_{(i)}}
  \;\in\; \bigl[0,\; 1 - \tfrac{1}{M}\bigr],
\end{equation}
attaining $0$ on a flat spectrum and $1 - 1/M$ when a single frequency carries everything.
The upper bound depends on $M$, so a Gini is comparable only across spectra of the same
length. This is one reason the random-orthogonal control below must preserve that length.

The participation ratio is computed on the energy, with
\begin{equation}
  \label{eq:pr}
  p_{k} \;=\; \frac{A_{k}^{2}}{\sum_{j} A_{j}^{2}},
  \qquad
  \mathrm{PR}(A) \;=\; \frac{1}{\sum_{k} p_{k}^{2}}
  \;=\; \frac{\bigl(\sum_{k} A_{k}^{2}\bigr)^{2}}{\sum_{k} A_{k}^{4}}
  \;\in\; [1, M] ,
\end{equation}
its standard definition. $\mathrm{PR} = M$ for a flat spectrum and $\mathrm{PR} = 1$ when one
frequency carries everything. It is an effective count of active frequencies, which is why it
belongs on energy: $p_{k}$ must be a probability over frequencies.

The two choices are not interchangeable and neither is a matter of taste. On one of our
$n = 113$ checkpoints, Gini on energy reads $0.902$ where Gini on amplitude reads $0.543$.
Feeding amplitude into the participation ratio gives $12.5$ where the energy convention gives
$4.31$. Each error was live in this project's code for months, and in the second
case it produced a failed prediction against a published value before it was found.

Three incompatible conventions appear in this literature, and their values are not
comparable. Writing $u$ for the spectrum being summarised,
\begin{equation}
  \label{eq:conventions}
  \underbrace{\mathrm{PR}(u) = \Bigl(\textstyle\sum_{k} p_{k}^{2}\Bigr)^{-1} \in [1, M]}_{\text{ours, and \citet{nguyen2026dlogclock}}}
  \qquad
  \underbrace{\mathrm{IPR}_{r}(u) = \bigl( \lVert u \rVert_{2r} / \lVert u \rVert_{2} \bigr)^{2r} \in (0, 1]}_{\text{\citet{doshi2024grokking}, with } r = 2}
\end{equation}
move in \emph{opposite} directions: ours decreases with sparsity, theirs increases, and
$\mathrm{IPR}_{2}(u) = 1/\mathrm{PR}(u)$ exactly when both are computed on the same
normalised energy. \citet{nanda2023progress} use a third. A sparsity number quoted without its
convention cannot be checked, and we state ours at every use.

\paragraph{Never one basis alone.} A purely additive signal scores $0.53$--$0.63$ in the
multiplicative basis, so a multiplicative Gini reported on its own is uninterpretable. Every
sparsity measurement in this paper is reported as a triple (multiplicative, additive, and a
random-orthogonal control), together with the count of key frequencies.

\subsection{Controls}
\label{sec:methods-controls}

The random-orthogonal control must rotate the residue axis, not the feature axis.

\begin{proposition}[Rotating features is a no-op]
\label{prop:noop}
Let $M \in \mathbb{R}^{G \times D}$ and let $Q \in \mathbb{R}^{D \times D}$ be orthogonal.
Then the amplitude spectrum \eqref{eq:amplitude} of $MQ$ equals that of $M$ exactly.
\end{proposition}

\begin{proof}
The transform \eqref{eq:transform} acts on the residue index, so
$\widehat{MQ}_{\mathbf{k}} = \widehat{M}_{\mathbf{k}} Q$. Since $Q$ is orthogonal it preserves
the Euclidean norm of each row, so
$\lVert \widehat{MQ}_{\mathbf{k}} \rVert = \lVert \widehat{M}_{\mathbf{k}} \rVert$ for every
$\mathbf{k}$, and $A_{\mathbf{k}}$ in \eqref{eq:amplitude} is a function of that norm alone.
\end{proof}

A control built that way does not look broken. It returns the signal exactly, which reads as
a result worth reporting (``the clock is no better than a random basis''). That is the
failure mode Proposition~\ref{prop:noop} exists to prevent. The control we use instead acts on
the \emph{residue} index,
\begin{equation}
  \label{eq:randorth}
  M \;\longmapsto\; Q^{*} M, \qquad Q \in \mathrm{U}(|G|) \text{ Haar-random},
\end{equation}
which preserves the spectrum length $M$ (so the Gini bound in \eqref{eq:gini} is unchanged)
and destroys the character structure. Both directions are asserted in the analysis code's
self-check.

Two further controls appear throughout. A \emph{shuffle} control permutes cells within the
object being measured, so it preserves its size and marginal distribution and destroys its
structure. A \emph{matched-failure} control is built only from runs that failed to
generalise. The last of these is worth stating as a rule. A binary accuracy gate answers ``is
this a data point''. It is the wrong instrument for ``is this a control''. One of our runs
ends at test accuracy $0.9779$, a model that has learned the circuit. Scored as an ungrokked
negative control, it read $74.8\%$ neuron tuning and destroyed the contrast the control
exists to provide. We therefore classify every run into three states, grokked,
near-grok and failed, and build controls only from the failed ones. This biases against our own
findings, which is the direction that is rarely checked.

\paragraph{Classifying a run.}
% STE descriptive
A run's endpoint is a window, never its last logged row. Under weight decay $1.0$, post-grok
accuracy excursions are common. Across $11$ long runs we observe $19$ excursions below $0.90$
after grokking. Every one is exactly one $200$-step logging interval wide, and $18$ of them
recover. The nineteenth merely begins at the last step at which an observation exists.

For the logged test accuracies $\mathrm{acc}_{1}, \dots, \mathrm{acc}_{S}$ of a run, we
define
\begin{equation}
  \label{eq:window}
  \bar{a} \;=\; \mathrm{median}\bigl( \mathrm{acc}_{S-9}, \dots, \mathrm{acc}_{S} \bigr),
  \qquad
  \text{state} =
  \begin{cases}
    \text{grokked}   & \bar{a} > 0.99, \\
    \text{near-grok} & 0.90 \le \bar{a} \le 0.99, \\
    \text{failed}    & \bar{a} < 0.90,
  \end{cases}
\end{equation}
and test accuracy is read by column name rather than by position, because our two
implementations write different quantities at the same index.

There are three states, not two: a binary gate turns every near-miss into whichever side it
falls on. The run above ends at $\bar{a} = 0.976$, so it is near-grok and not a
control. Controls are built only from runs in the \emph{failed} state. The threshold $0.99$ is
the grokking criterion of every training script. The threshold $0.90$ is the floor of the
excursion count above, and it was fixed before any control arm was scored.
% STE end

\subsection{The ablation protocol}
\label{sec:methods-causal}

Sparsity says a representation is concentrated. It does not say the concentration is used.

Let $L \in \mathbb{R}^{G \times G \times n}$ be the logit tensor over a grid of units,
re-indexed into character coordinates by \eqref{eq:transform} on both input axes, so that
$\widehat{L}_{\mathbf{k}, \mathbf{k}'}$ is its component at the character pair
$(\mathbf{k}, \mathbf{k}')$. For a set $S \subseteq \widehat{G} \times \widehat{G}$ define the
spectral projection $\Pi_{S}$ by
\begin{equation}
  \label{eq:project}
  \widehat{\Pi_{S} L}_{\mathbf{k}, \mathbf{k}'} \;=\;
  \begin{cases}
    \widehat{L}_{\mathbf{k}, \mathbf{k}'} & (\mathbf{k}, \mathbf{k}') \in S, \\
    0 & \text{otherwise,}
  \end{cases}
\end{equation}
and write $\mathcal{L}(\cdot)$ for cross-entropy against the true targets. The three
quantities we report are then
\begin{equation}
  \label{eq:ablation}
  \underbrace{\mathcal{L}(L)}_{\text{baseline}}
  \qquad
  \underbrace{\mathcal{L}(\Pi_{K} L)}_{\text{restricted}}
  \qquad
  \underbrace{\mathcal{L}(\Pi_{K^{c}} L)}_{\text{excluded}}
\end{equation}
for a key set $K$, where $K^{c}$ is its complement. \emph{Restricted} keeps only $K$;
\emph{excluded} removes exactly $K$ and keeps everything else. We report all three as absolute
losses side by side, following \citet{chughtai2023toy}.

\paragraph{What this intervention acts on.} $\Pi_{S}$ acts on the
\emph{logit tensor} of a trained model. The forward pass is not re-run, no weight is ablated
and no activation is patched. This is \citet{nanda2023progress}'s restricted and excluded
loss protocol, applied in a coordinate they did not need, and what it licenses is a statement
about the output: the logits are carried by $K$ and are destroyed when exactly $K$ is removed.
It is not a statement about which internal component computes them.
\citet{chen2026multiplication} name ablations \emph{and} activation patching as the evidence
their account lacks; this protocol supplies the first. The title, the abstract and the heading
of Section~\ref{sec:causal} therefore name the ablation rather than assert causation, and
where the claim ledger of Appendix~\ref{app:ledger} uses the word \emph{causal} it means this
protocol and no other.

\paragraph{The one place we do re-run the forward pass.}
% STE descriptive
At $n = 113$ and $n = 121$ only, the same projection $\Pi_{S}$ of \eqref{eq:project} is
applied to the \emph{embedding} $W_{E}$ rather than to the logit tensor. The edit is made in
the multiplicative character basis. The network then runs forward unmodified
(Section~\ref{sec:causal-internal}). Softmax attention and the ReLU MLP run on top of the edited weights. So, unlike a mask on the output, nothing here is
algebraically forced. Non-unit rows and the read-out token row carry no character index and
pass through untouched.

The experiment was pre-registered before the models it uses were retrained. Its three
thresholds are inherited from criteria this paper fixes elsewhere for other purposes. They are
the $100\times$ of Section~\ref{sec:strata}, the $0.90$ generalisation floor, and
$p < 0.01$. So none was chosen after a number was seen.
% STE end

The two nonlinearities lie downstream of the edit, so the result is not a rearrangement of
the output, and the scoping of the logit-tensor protocol does not apply to it. It applies
everywhere else in this paper. We keep the two apart deliberately: two moduli with a weight
intervention do not retroactively upgrade the other twenty-one, measured a different way.

% STE descriptive
\begin{figure}[t]
  \centering
  \includegraphics[width=\textwidth]{../figures/F3_ablation_scheme.pdf}
  \caption[The two projections, drawn on the plane they act on.]{\textbf{The two projections, drawn on the plane they act on.} Each panel is the
    character plane of the logits at $n = 121$ seed $0$, in discrete-log coordinates, with
    DC at the centre. Shade is $\log_{10}$ of the magnitude
    $\lVert \widehat{L}_{\kappa_{a}, \kappa_{b}} \rVert$ summed over the output axis.
    \emph{Left:} the model as trained. The energy that matters lies on the diagonal
    $\Delta$, which is where \eqref{eq:diagonal} says a character circuit must put it.
    \emph{Centre:} $\Pi_{K}$ keeps the $4$ key characters and their conjugates and zeroes
    the other $12{,}091$ pairs.\par
    Rings mark the survivors, because nine kept pixels out of $110^{2}$ would otherwise read
    as nothing there. \emph{Right:} $\Pi_{K^{c}}$ removes exactly those nine and keeps the
    rest. The cross-entropy under each panel is printed beneath it: restricted is at
    or below baseline, and excluded is seven orders of magnitude worse. So the key characters
    are both sufficient and necessary on this run. The figure regenerates from
    \texttt{src/viz/plots.py::ablation\_scheme}. The spectrum, the key set and all three
    losses come from \texttt{src/analysis/ablation.py}, the same code
    Section~\ref{sec:causal} reports.}
  \label{fig:ablation-scheme}
\end{figure}
% STE end

\paragraph{The statistic is a permutation $p$, not a ratio to a control mean.} Drawing
$R_{1}, \dots, R_{B}$ uniformly at random from the character sets of cardinality $|K|$, with
$B = 10{,}000$ and $b$ the number of draws at least as damaging as the key set,
\begin{equation}
  \label{eq:permp}
  \hat{p} \;=\; \frac{1 + b}{B + 1},
  \qquad
  b \;=\; \sum_{b'=1}^{B}
  \mathbf{1}\Bigl[\, \mathcal{L}(\Pi_{R_{b'}^{c}} L) \;\ge\; \mathcal{L}(\Pi_{K^{c}} L) \,\Bigr],
\end{equation}
the fraction of equal-sized random sets at least as damaging as the key set.
The $(1+b)/(B+1)$ form \citep{phipson2010permutation} is deliberate, and the naive $b/B$ is a trap we fell into. The
ratio $b/B$ reads exactly $0$ whenever no draw is as damaging as the key set, and a Monte-Carlo
$p$ of $0$ claims a resolution the design does not have. An earlier draft of this paper
printed $p = 0.0000$ to four decimals from a $200$-draw null whose smallest expressible
value is $1/201 \approx 0.005$. Here the floor is $\hat{p} = 1/10{,}001 \approx 1.0 \times
10^{-4}$, and where $b = 0$ the text states it in words as well as in the estimate. It is
threshold-free and scale-free, and it is the statistic every dependence claim in this paper
rests on. We previously reported
$\text{excluded} / \text{mean}(\text{control})$ over $20$ draws and have retired it. The
control distribution is bimodal: roughly $73\%$ of random character removals leave the loss
untouched, because the model does not use those characters, and the remainder are
catastrophic. Its mean is therefore decided by how many catastrophic draws happen to land, and
at $n = 121$ seed $0$ that ratio ranges over $21\times$ to $440\times$ across control random
seeds alone, on identical data. Where a descriptive ratio is useful we quote the median
control. The same permutation logic, at $20{,}000$ draws, is used for the enrichment test in
Section~\ref{sec:crt-law}.

\paragraph{Multiplicity.} This paper's permutation tests form a
family of $320$: $157$ measurable stratum-measurements in the primary arm and $26$ in the
engine arm, plus $137$ single-modulus ablations across eight sweeps and two runs of a retired engine arm. They test one
hypothesis repeatedly rather than searching for a significant one among many, but the count
is what it is and the reader should not have to reconstruct it. At $B = 10{,}000$ the floor
$1.0 \times 10^{-4}$ clears a Benjamini--Hochberg correction at $\alpha = 0.01$ and a
Bonferroni correction at $\alpha = 0.05$ (per-test $1.6 \times 10^{-4}$). It does not
clear Bonferroni at $\alpha = 0.01$, which would require a per-test level of
$3.1 \times 10^{-5}$ and hence $B \ge 32{,}000$. That comparison fails, and it is stated here
rather than omitted. The draw count was chosen for cost: $B = 10{,}000$ is about ten hours of
local CPU across the family and $B = 32{,}000$ about thirty. This reason is recorded here so
it is not later mistaken for a statistical judgement.

\paragraph{Two guards against circularity in the ablated set.}
For the unit-stratum results the key set is recovered independently of the ablation, by
embedding norm, and the set recovered by ablation rank agrees with it. For the
per-stratum results of Section~\ref{sec:strata} no key-frequency detector is used at
all: the ablated set is the \emph{whole diagonal}, fixed in advance by
Theorem~\ref{thm:stratum}. Selecting a set by ablation rank and then permutation-testing that
same set would be circular, and the second design excludes it by construction.

\paragraph{Three further statistics.} The \emph{enrichment} of a predicted
frequency set $P$ against a spectrum $A$ compares its mean mass to the global mean,
\begin{equation}
  \label{eq:enrich}
  \mathrm{enr}(P) \;=\;
  \frac{\frac{1}{|P|}\sum_{k \in P} A_{k}}{\frac{1}{M}\sum_{k} A_{k}} ,
\end{equation}
and is scored by the same permutation logic as \eqref{eq:permp} over $20{,}000$ random sets of
cardinality $|P|$. The \emph{diagonal energy share} of a stratum measures how much of the
logit energy sits on $\Delta$ of \eqref{eq:diagonal}, against the share an unstructured tensor
would give,
\begin{equation}
  \label{eq:dshare}
  D \;=\; \frac{\sum_{(\mathbf{k},\mathbf{k}') \in \Delta} \lVert \widehat{L}_{\mathbf{k},\mathbf{k}'} \rVert^{2}}
               {\sum_{\mathbf{k},\mathbf{k}'} \lVert \widehat{L}_{\mathbf{k},\mathbf{k}'} \rVert^{2}} ,
  \qquad
  R \;=\; \frac{D}{1 / (|G_{m}| + 1)} ,
\end{equation}
so $R$ is the ratio to exact chance. Finally a neuron $h : G \times G \to \mathbb{R}$ is
\emph{tuned} to frequency $\mathbf{k}$ when a single frequency's eight-bin family explains
most of its DC-free two-dimensional spectral energy,
\begin{equation}
  \label{eq:tuned}
  \mathrm{tune}(h) \;=\; \max_{\mathbf{k} \neq 0}
  \frac{\sum_{(\mathbf{j},\mathbf{j}') \in F(\mathbf{k})} \lVert \widehat{h}_{\mathbf{j},\mathbf{j}'} \rVert^{2}}
       {\sum_{(\mathbf{j},\mathbf{j}') \neq 0} \lVert \widehat{h}_{\mathbf{j},\mathbf{j}'} \rVert^{2}}
  \;>\; 0.85 ,
\end{equation}
where $F(\mathbf{k})$ is the eight-bin family $\{(\pm\mathbf{k}, 0), (0, \pm\mathbf{k}),
(\pm\mathbf{k}, \pm\mathbf{k})\}$. The threshold $0.85$ and the eight-bin family are
\citet{nanda2023progress}'s; scoring a single bin instead reads $9.6\%$ where the correct
statistic reads $85.0\%$ on the same model.

\paragraph{The two-dimensional bin convention, which is load-bearing.} On a grid indexed so
that multiplication is addition, a function of the product concentrates on the diagonal and a
function of the ratio on the antidiagonal:
\begin{equation}
  \label{eq:bins}
  f(a + b) \;\longleftrightarrow\; (k, k),
  \qquad
  f(a - b) \;\longleftrightarrow\; (k, -k).
\end{equation}
Reported the other way round, \eqref{eq:bins} inverts every conclusion drawn from it. We
verified it by planting signals of known frequency and locating them, never by deriving it.
Three separate bin-convention errors in this project were each found that way and none by
reasoning.

\paragraph{A generation diagnostic.} A key set that does not \emph{generate}
the character group cannot implement an exact clock, since two units on which every key
character agrees cannot be separated by any weighting of those characters. We report this as a
per-run diagnostic, and every key set recovered here generates its full group.

\subsection{What the measurements are worth}
\label{sec:methods-validity}

\paragraph{Noise in a single-checkpoint Gini.} Measured on three moduli at three
seeds over the final $12{,}000$ steps, $\mathrm{Gini}_{\text{mult}}$ has within-run standard
deviation $0.02$ to $0.044$. At $n = 125$ seed $2$ it wanders over $0.498$--$0.621$, a range of
$0.122$, with no direction. This is oscillation about a flat mean and not a trend. The
statistic plateaus by roughly $16{,}000$ to $20{,}000$ steps and then jitters. Against the
$0.10$ threshold we use for ``the clock survives'', a single checkpoint is therefore about a
$2\sigma$ test, and only about $1.6\sigma$ at $n = 125$. Where the comparison matters we
average over the final $10{,}000$ steps and report the within-run spread. It follows that part
of the across-seed spread reported for a single-checkpoint measurement is estimator noise
rather than seed variance. The two were not separated in the earlier measurements, and they are
not re-attributed here.

\paragraph{Convergence at forty thousand steps.} For a
statistic $\sigma$ sampled along a run at steps $t$, write its \emph{drift} against the final
checkpoint $T$ as
\begin{equation}
  \label{eq:drift}
  \mathrm{drift}(\sigma) \;=\;
  \frac{1}{|W|}\sum_{t \in W} \frac{\bigl|\sigma(t) - \sigma(T)\bigr|}{|\sigma(T)|},
  \qquad W = \text{the final } 10{,}000 \text{ steps},
\end{equation}
and call a seed converged when $\mathrm{drift} < 0.05$. Requiring that in at least four of
five seeds, $n = 113$
converges ($2.3\%$ participation-ratio drift, $5/5$ seeds) and $n = 121$ converges ($2.8\%$,
$4/5$). Four moduli do not: $n = 125$ at $5.6\%$ ($2/5$ seeds), $n = 119$ at $5.9\%$ ($3/5$),
$n = 120$ at $7.7\%$ ($3/5$) and $n = 165$ at $4.8\%$ ($3/5$). Every participation ratio at
those four is measured at a non-stationary point. Rather than re-measure at a longer horizon,
we retire the participation ratio at composite moduli and do not use it for cross-modulus
comparison there.
That a representation statistic need not be monotone in training is not a peculiarity of our
setup. \citet{demoss2025complexity} track a compression-based complexity across grokking and
find it \emph{rises} through memorisation and only then falls. The value read at any single
checkpoint therefore depends on where in that arc the checkpoint sits. Our drift criterion is the
narrow, per-statistic version of that caution, and the moduli it fails on are exactly the ones
we decline to compare.

\paragraph{The fibre residual.} Consequence~2 of Theorem~\ref{thm:stratum} says the target is
constant on the fibres of \eqref{eq:fibration}. To test whether the \emph{logits} are too, we
average a block $B$ over its fibres, write $\bar{B}$ for the result lifted back to the block,
and report the share of block variance that averaging destroys:
\begin{equation}
  \label{eq:resid}
  \mathrm{resid}(B) \;=\;
  \frac{\bigl\lVert B - \bar{B} \bigr\rVert_{F}^{2}}
       {\bigl\lVert B - \mathrm{mean}(B) \bigr\rVert_{F}^{2}} \;\in\; [0, 1].
\end{equation}
Averaging always smooths, so a small residual on its own means nothing. We therefore compare
against $\mathrm{resid}$ computed on
\emph{permuted} fibres of the same sizes, which have no algebraic meaning. This is the
size-matched defence that three retracted claims of ours needed and did not have.

\paragraph{Effect sizes for a pre-registered two-cell comparison.} Where an experiment asks
what fraction of a known gap a variable explains, we fix the two endpoints in the
pre-registration and report
\begin{equation}
  \label{eq:effectsize}
  f \;=\; \frac{\sigma_{\text{observed}} - \sigma_{\text{baseline}}}
                {\sigma_{\text{target}} - \sigma_{\text{baseline}}},
\end{equation}
with the decision rule committed before the observation exists: held at $f \ge 0.70$, not
held at $f \le 0.30$, and \emph{partial}, reported as such, in between.
Because the endpoints are fixed in advance, $f$ can exceed $1$ or fall below $0$, and we
report the value rather than clipping it.

\paragraph{Training precision changes an absolute sparsity number in PyTorch.} The same
PyTorch code at $n = 113$ reads $\mathrm{Gini}_{\text{mult}} = 0.5791$ at float32 and $0.7728$
at float64, while our from-scratch engine reads $0.7671$ and $0.7536$ at the two dtypes and is
effectively unaffected. The effect is an interaction between implementation and precision, not
a main effect of precision. We tested and retracted the main-effect claim, and
Appendix~\ref{app:retractions} records that. Every sweep in this paper is float32, and we
name the regime beside every absolute sparsity number. Comparisons between moduli, and
comparisons within one arm, are unaffected, which is what almost every result here rests on.
